\chapter{Methodology}
\label{chap:methodology}
\label{methodology}
The proposed system follows a two-stage pipeline. In the first stage, the LMM region is localized from the ultrasound image. In the second stage, the localized image is classified into one of the four Heckmatt grades.

The ROI extraction stage was first studied using the LUMINOUS dataset, where segmentation masks were converted into bounding boxes. The ROI detector was then used to crop the muscle region from the clinical GMC dataset. Earlier experiments used U-Net and SAM for segmentation, but the final practical solution used YOLOv8 because it can directly predict a bounding box at inference time without requiring a prompt.

For classification, the cropped or full ultrasound images were passed to transfer learning based models. Several backbones were tested, including EfficientNet-B0, ResNet, ViT-B16, SigLIP, and the MedGemma vision encoder. ResNet was included as a stable convolutional baseline because residual learning is well suited for training deeper networks without the same optimization problems as plain CNNs \cite{he2016resnet}. ViT was used because the texture pattern in ultrasound is not always local, and the self-attention based model can compare image regions more globally.

The first classification setup was a direct four-class experiment. Each image was assigned to H1, H2, H3, or H4, and the model was trained with the usual train, validation, and test split. Because the class counts were uneven, accuracy alone was not enough. Balanced accuracy and confusion matrices were used to see whether the model was only learning the dominant classes. Weighted loss was also tested to reduce the bias caused by the larger H2 group, following the general idea of imbalanced learning \cite{HeGarcia2009,Cui2019}.

Preprocessing was tested because the ultrasound images were not visually uniform. Some images had strong speckle and some had lower contrast. Gaussian blur, grayscale conversion, ImageNet normalization, and augmentation were tried in different runs. The aim was not to remove all texture, since texture is part of the Heckmatt signal, but to reduce unwanted noise and make the model less sensitive to small acquisition differences. This part of the methodology is related to ultrasound despeckling work \cite{YuActon2002,Loizou2005} and image augmentation literature \cite{Shorten2019,Cubuk2019}.

After the direct four-class runs, pairwise binary classifiers were trained for all class pairs: H1 vs H2, H1 vs H3, H1 vs H4, H2 vs H3, H2 vs H4, and H3 vs H4. This was done to check which boundaries were actually separable. If a pairwise model performs well for distant grades but poorly for adjacent grades, then the problem is not only the final classifier. It also means the visual boundary between those grades is weak or inconsistent in the available data.

A hierarchical method was also tested because Heckmatt grades have an order. Two directions were used. The low-to-high cascade first checked for H1, then H2, and finally separated H3 from H4. The high-to-low cascade did the reverse by first checking H4, then H3, and finally separating H2 from H1. This made it possible to see whether the model worked better when severe cases were separated first or when normal cases were separated first.

The most recent audit phase used 606 labelled Heckmatt images from the GMC and LUMINOUS image sets. For the audit, H1 and H2 were mapped to the normal class, while H3 and H4 were mapped to the abnormal class. A ViT model was used to compare predicted normal/abnormal labels against these assumptions. A disagreement was counted when the model prediction did not match the assumed label, and a high-confidence disagreement was counted when the prediction confidence was at least 0.8. This was not used as a final medical decision. It was used as a practical way to find images that should be reviewed again by the doctor.

The project also produced an annotation workflow for the relabeling stage. Images can be uploaded into the tool, shown to the doctor one by one, assigned a corrected Heckmatt label, and marked with a bounding box around the muscle region. At the end, the corrected labels and bounding-box data can be downloaded directly. This matters because the audit showed that some labels are likely inconsistent, and the doctor is currently correcting those cases.


\begin{figure}[H]
\centering
\safeincludegraphics[width=0.4\textwidth]{Screenshot 2026-02-25 232024.png}
\caption{YOLO fine-tuning pipeline}
\label{fig:yolo_finetune}
\end{figure}

\clearpage

\begin{figure}[H]
\centering
\safeincludegraphics[width=0.7\textwidth]{datasetDistribution.png}
\caption{Dataset Distribution by Grades}
\label{fig:dataset_distribution}
\end{figure}



\begin{figure}[H]
\centering
\safeincludegraphics[width=0.9\textwidth]{Architecture-of-the-YOLO-network.png}
\caption{YOLO architecture. Source: ResearchGate}
\label{fig:yolo_arch}
\end{figure}

\begin{figure}[H]
\centering
\safeincludegraphics[width=0.9\textwidth]{yolo_arch.png}
\caption{Classification methodology}
\label{fig:classification_methodology}
\end{figure}

\begin{figure}[H]
\centering
\safeincludegraphics[width=0.7\textwidth]{EfficientNet-Architecture.png}
\caption{EfficientNet architecture}
\label{fig:efficientnet_arch}
\end{figure}


\begin{figure}[H]
\centering
\safeincludegraphics[width=0.9\textwidth]{hierarchical-inf-pipe.png}
\caption{Hierarchical inference architecture}
\label{fig:hierarchical_arch}
\end{figure}

\begin{figure}[H]
\centering
\safeincludegraphics[width=0.9\textwidth]{label-audit.png}
\caption{Binary label audit workflow}
\label{fig:audit_arch}
\end{figure}

\section{Evaluation Metrics}
\label{metrics}
For segmentation and ROI localization experiments, the following metrics were used:

\textbf{Intersection over Union (IoU) or Jaccard Index:}
\[
\text{IoU} = \frac{|A \cap B|}{|A \cup B|}
\]

\textbf{Dice Coefficient:}
\[
\text{Dice} = \frac{2|A \cap B|}{|A| + |B|}
\]

\textbf{Pixel Accuracy:}
\[
\text{Pixel Accuracy} = \frac{\text{Correctly classified pixels}}{\text{Total pixels}}
\]

For classification, the following metrics were used:

\textbf{Accuracy:}
\[
\text{Accuracy} = \frac{TP + TN}{TP + TN + FP + FN}
\]

\textbf{Precision:}
\[
\text{Precision} = \frac{TP}{TP + FP}
\]

\textbf{Recall / Sensitivity:}
\[
\text{Recall} = \frac{TP}{TP + FN}
\]

\textbf{Specificity:}
\[
\text{Specificity} = \frac{TN}{TN + FP}
\]

\textbf{F1 Score:}
\[
\text{F1 Score} = \frac{2 \cdot \text{Precision} \cdot \text{Recall}}{\text{Precision} + \text{Recall}}
\]

\textbf{Balanced Accuracy:}
\[
\text{Balanced Accuracy} = \frac{1}{C}\sum_{i=1}^{C}\text{Recall}_{i}
\]

Balanced accuracy was especially important in this project because the Heckmatt classes were not equally represented.

